\section{Von Neumann Model and ISA}

\begin{definition}
The \textbf{Von Neumann model}, proposed in 1946, defines a universal architecture for general-purpose computers consisting of five fundamental components:
\begin{itemize}
    \item \textbf{Memory:} Stores both the program (instructions) and the data.
    \item \textbf{Processing Unit:} Performs arithmetic and logical operations. It contains the Arithmetic Logic Unit (ALU) and small temporary storage (registers).
    \item \textbf{Control Unit:} Orchestrates the execution of instructions. It contains the Program Counter (PC), which tracks the address of the next instruction, and the Instruction Register (IR), which holds the current instruction.
    \item \textbf{Input:} Devices to get information into the computer (e.g., keyboard).
    \item \textbf{Output:} Devices to get information out (e.g., monitor).
\end{itemize}
\end{definition}

\begin{important}
Modern computing is based on two key properties of this model:
\begin{itemize}
    \item \textbf{Stored Program:} Instructions are stored in a linear memory array alongside data. The hardware interprets a memory value as an instruction or data depending on which phase of the cycle it is in.
    \item \textbf{Sequential Processing:} Instructions are processed one at a time in sequence, unless a control flow instruction explicitly changes the Program Counter.
\end{itemize}
\end{important}

\begin{remark}
The process of executing a single instruction follows a specific sequence of phases:
\begin{itemize}
    \item 1. \textbf{FETCH:} Obtain instruction from memory into the IR and increment the PC.
    \item 2. \textbf{DECODE:} Identify the opcode and determine the control signals needed.
    \item 3. \textbf{EVALUATE ADDRESS:} Calculate the memory address of operands if needed (e.g., for loads).
    \item 4. \textbf{FETCH OPERANDS:} Get the required source values from memory or registers.
    \item 5. \textbf{EXECUTE:} Perform the actual operation in the ALU.
    \item 6. \textbf{STORE RESULT:} Write the final value into the destination register or memory.
\end{itemize}
\end{remark}

\begin{definition}
\textbf{Address Space:} The total number of uniquely identifiable memory locations. (LC-3: $2^{16}$; MIPS: $2^{32}$).
\end{definition}

\begin{definition}
\textbf{Addressability:} The number of bits stored at each address (e.g., byte-addressable (8 bit) vs. word-addressable (32 bit)).
\end{definition}

\begin{definition}
\textbf{Endianness:} The order of bytes in a word. \textbf{Big Endian} stores the Most Significant Byte (MSB) at the lowest address, while \textbf{Little Endian} stores the Least Significant Byte (LSB) at the lowest address.
\end{definition}

\begin{example}
To read from memory:
\begin{enumerate}
    \item Load the \textbf{Memory Address Register (MAR)} with the target address.
    \item The memory hardware places the data from that address into the \textbf{Memory Data Register (MDR)}.
\end{enumerate}
To write, both MAR (address) and MDR (data) are loaded before activating the write-enable signal.
\end{example}

\subsection{Instruction Set Architecture (ISA)}

\begin{definition}
The Instruction Set Architecture (ISA) serves as the "contract" or interface between the software (compiler/programmer) and the hardware (microarchitecture). It specifies:
\begin{itemize}
    \item \textbf{Memory Organization:} Address space and addressability.
    \item \textbf{Register Set:} Number, size, and usage of registers (e.g., 8 GPRs in LC-3, 32 in MIPS).
    \item \textbf{Instruction Set:} Opcodes, data types, and addressing modes.
\end{itemize}
\end{definition}

\begin{important}
Instructions are encoded as binary "words" containing:
\begin{itemize}
    \item \textbf{Opcode:} Specifies the operation to perform (e.g., ADD, LD).
    \item \textbf{Operands:} Specify where to find the source data and where to put the result.
\end{itemize}
\end{important}

\begin{lemma}
The semantic gap refers to the distance between high-level language constructs and the ISA. 
\begin{itemize}
    \item \textbf{CISC (Complex Instruction Set Computing):} Small semantic gap; instructions do a lot of work (e.g., x86).
    \item \textbf{RISC (Reduced Instruction Set Computing):} Large semantic gap; simple instructions serve as primitives (e.g., MIPS, LC-3).
\end{itemize}
\end{lemma}

\begin{definition}
Mechanisms used to specify where an operand is located:
\begin{itemize}
    \item \textbf{Immediate:} The operand is a constant within the instruction.
    \item \textbf{Register:} The operand is located in a general-purpose register.
    \item \textbf{PC-Relative:} The address is calculated as $PC + offset$.
    \item \textbf{Indirect:} The instruction points to a memory location that contains the \textit{actual} address of the operand.
    \item \textbf{Base + Offset:} The address is calculated as $Register + offset$.
\end{itemize}
\end{definition}

\subsection{Assembly}

\begin{definition}
\textbf{Machine code} is the binary representation (0s and 1s) processed by hardware. \textbf{Assembly language} is a human-readable representation using mnemonics (e.g., `ADD`) and labels for memory locations.
\end{definition}

\begin{example}
High-level: $a = b + c$
\begin{itemize}
    \item \textbf{LC-3 Assembly:} `ADD R0, R1, R2` (R0 is destination).
    \item \textbf{MIPS Assembly:} `add \$s0, \$s1, \$s2` (\$s0 is destination).
\end{itemize}
\end{example}

\begin{important}
In RISC (Reduced Instruction Set Computing) architectures, operands must be moved into registers before processing.
\begin{itemize}
    \item \textbf{Load (LD/LDR/lw):} Reads data from memory into a register.
    \item \textbf{Store (ST/STR/sw):} Writes data from a register into memory.
\end{itemize}
\end{important}

\begin{remark}
The sequential execution order can be altered using:
\begin{itemize}
    \item \textbf{Conditional Branches:} Executed only if a condition is met. LC-3 uses \textbf{Condition Codes} (N, Z, P) set by previous instructions. MIPS often compares two registers directly (e.g., `beq $s0, $s1, Label`).
    \item \textbf{Unconditional Jumps:} Always change the PC (e.g., `JMP` in LC-3, `j` in MIPS).
\end{itemize}
\end{remark}